\QHeading{第1题　口算。}
\par\noindent


\begin{oralgrid}
\OralItem{1}{2.45\times10}{24.5}
\OralItem{2}{0.138\times100}{13.8}
\OralItem{3}{0.009\times1000}{9}
\OralItem{4}{2.5\times1000}{2500}
\OralItem{5}{35.8\div10}{3.58}
\OralItem{6}{0.4\div100}{0.004}
\OralItem{7}{30\div100}{0.3}
\OralItem{8}{5\div1000}{0.005}\end{oralgrid}
\par

\QHeading{口算思路补充（对应上表题号）}

\Answer{上面列出得数，下面按小题给出不同计算思路。计数单位法适用于小数，整数题可比较凑整、拆分或商不变等方法。}
\MethodTitle{第7小题：$30\div 100$}

\MethodBlock{口算思路一：同时去掉相同个数的0}{}{\begin{aligned}[t]&\phantom{={}}30\div 100\\
&=(30\div 10)\div (100\div 10)\\
&=3\div 10\\
&=0.3\end{aligned}}


\par\noindent

\QHeading{第2题　求近似数。}

\PlainItem{精确到十分位：$1.07\approx${\color{answerred}$1.1$}}
\par

\Equation{0.53\approx\textcolor{answerred}{0.5}}
\Equation{0.98\approx\textcolor{answerred}{1.0}}
\Equation{10.094\approx\textcolor{answerred}{10.1}}
\PlainItem{保留整数：$67.4\approx${\color{answerred}$67$}}
\Equation{0.83\approx\textcolor{answerred}{1}}
\Equation{9.89\approx\textcolor{answerred}{10}}
\Equation{4.8\approx\textcolor{answerred}{5}}
\par\noindent

\QHeading{第3题　用简便方法计算下面各题。}

\CalcBlock[方法一：简便计算]{1}{\begin{aligned}[t]&\phantom{={}}\color{black}357+288+143\\&\color{answerred}=357+143+288\\&\color{answerred}=(357+143)+288\\&\color{answerred}=500+288\\&\color{answerred}=788\end{aligned}}
\par
\MethodBlock{方法二：按运算顺序直接计算（对照）}{本法只作结果核对；本题要求简便计算，前面的简便解法仍保留。}{\begin{aligned}[t]&\phantom{={}}357+288+143\\
&=645+143\\
&=788\end{aligned}}


\CalcBlock[方法一：简便计算]{2}{\begin{aligned}[t]&\phantom{={}}\color{black}58+39+42+61\\&\color{answerred}=58+42+39+61\\&\color{answerred}=(58+42)+(39+61)\\&\color{answerred}=100+100\\&\color{answerred}=200\end{aligned}}
\MethodBlock{方法二：交换、结合后先凑整}{保留每个数原来的加减号，先把能凑整的部分放在一起。}{\begin{aligned}[t]&\phantom{={}}58+39+42+61\\
&=(58+42)+39+61\\
&=100+39+61\\
&=139+61\\
&=200\end{aligned}}

\MethodBlock{方法三：按运算顺序直接计算（对照）}{本法只作结果核对；本题要求简便计算，前面的简便解法仍保留。}{\begin{aligned}[t]&\phantom{={}}58+39+42+61\\
&=97+42+61\\
&=139+61\\
&=200\end{aligned}}

\CalcBlock[方法一：简便计算]{3}{\begin{aligned}[t]&\phantom{={}}\color{black}3065-738-1065\\&\color{answerred}=3065-1065-738\\&\color{answerred}=2000-738\\&\color{answerred}=1262\end{aligned}}
\MethodBlock{方法二：先合并要减去的部分}{连续减去几个数，可以先把这些数相加再减。}{\begin{aligned}[t]&\phantom{={}}3065-738-1065\\
&=3065-(738+1065)\\
&=3065-1803\\
&=1262\end{aligned}}

\MethodBlock{方法三：按运算顺序直接计算（对照）}{本法只作结果核对；本题要求简便计算，前面的简便解法仍保留。}{\begin{aligned}[t]&\phantom{={}}3065-738-1065\\
&=2327-1065\\
&=1262\end{aligned}}

\CalcBlock[方法一：简便计算]{4}{\begin{aligned}[t]&\phantom{={}}\color{black}7755-(2187+755)\\&\color{answerred}=7755-2187-755\\&\color{answerred}=7755-755-2187\\&\color{answerred}=7000-2187\\&\color{answerred}=4813\end{aligned}}
\MethodBlock{方法二：去括号后计算}{加号后的括号去掉不变号，减号后的括号去掉要改变各项的加减号。}{\begin{aligned}[t]&\phantom{={}}7755-(2187+755)\\
&=7755-2187-755\\
&=5568-755\\
&=4813\end{aligned}}

\MethodBlock{方法三：按运算顺序直接计算（对照）}{本法只作结果核对；本题要求简便计算，前面的简便解法仍保留。}{\begin{aligned}[t]&\phantom{={}}7755-(2187+755)\\
&=7755-2942\\
&=4813\end{aligned}}

\CalcBlock[方法一：简便计算]{5}{\begin{aligned}[t]&\phantom{={}}\color{black}497-299\\&\color{answerred}=497-(300-1)\\&\color{answerred}=497-300+1\\&\color{answerred}=197+1\\&\color{answerred}=198\end{aligned}}
\MethodBlock{方法二：补差法}{从减数先补到一个整数，再补到被减数，两段差相加。}{\begin{aligned}[t]&\phantom{={}}497-299\\
&=(300-299)+(497-300)\\
&=1+197\\
&=198\end{aligned}}

\CalcBlock{6}{\begin{aligned}[t]&\phantom{={}}\color{black}278+98\\&\color{answerred}=278+(100-2)\\&\color{answerred}=278+100-2\\&\color{answerred}=378-2\\&\color{answerred}=376\end{aligned}}
\CalcBlock[方法一：简便计算]{7}{\begin{aligned}[t]&\phantom{={}}\color{black}263-52-163\\&\color{answerred}=263-163-52\\&\color{answerred}=100-52\\&\color{answerred}=48\end{aligned}}
\MethodBlock{方法二：先合并要减去的部分}{连续减去几个数，可以先把这些数相加再减。}{\begin{aligned}[t]&\phantom{={}}263-52-163\\
&=263-(52+163)\\
&=263-215\\
&=48\end{aligned}}

\MethodBlock{方法三：按运算顺序直接计算（对照）}{本法只作结果核对；本题要求简便计算，前面的简便解法仍保留。}{\begin{aligned}[t]&\phantom{={}}263-52-163\\
&=211-163\\
&=48\end{aligned}}

\CalcBlock[方法一：简便计算]{8}{\begin{aligned}[t]&\phantom{={}}\color{black}257-168-32\\&\color{answerred}=257-(168+32)\\&\color{answerred}=257-200\\&\color{answerred}=57\end{aligned}}
\MethodBlock{方法二：按运算顺序直接计算（对照）}{本法只作结果核对；本题要求简便计算，前面的简便解法仍保留。}{\begin{aligned}[t]&\phantom{={}}257-168-32\\
&=89-32\\
&=57\end{aligned}}

\CalcBlock[方法一：简便计算]{9}{\begin{aligned}[t]&\phantom{={}}\color{black}607-102\\&\color{answerred}=607-(100+2)\\&\color{answerred}=607-100-2\\&\color{answerred}=507-2\\&\color{answerred}=505\end{aligned}}
\MethodBlock{方法二：补差法}{从减数先补到一个整数，再补到被减数，两段差相加。}{\begin{aligned}[t]&\phantom{={}}607-102\\
&=(110-102)+(607-110)\\
&=8+497\\
&=505\end{aligned}}

\par\noindent

\QHeading{第4题　解决问题。}

\Problem{文化路小学准备发放500份宣传材料，前4天平均每天发放80份，剩下的在3天内发完。后3天平均每天要发放多少份？}
\par


\MethodTitle{方法一}
\Answer{{\Formula{\begin{aligned}[t]&\phantom{={}}(500-80\times4)\div3\\&=(500-320)\div3\\&=180\div3\\&=60\text{（份）}\end{aligned}}}\par 答：后3天平均每天要发60份。}

\MethodTitle{方法二：用乘法检验后3天每天的发放量}
\Answer{前4天共发320份，剩180份；若每天60份，3天正好180份。}
\Answer{\textbf{前4天发放数}\par \Formula{\begin{aligned}[t]&\phantom{={}}80\times 4\\
&=320\text{（份）}\end{aligned}}}
\Answer{\textbf{剩余数}\par \Formula{\begin{aligned}[t]&\phantom{={}}500-320\\
&=180\text{（份）}\end{aligned}}}
\Answer{\textbf{每天60份，3天合计}\par \Formula{\begin{aligned}[t]&\phantom{={}}60\times 3\\
&=180\text{（份）}\end{aligned}}}
\Answer{答：后3天平均每天要发60份。}
